% TikZ 图：动能定理示意
% Day 08：功、功率与动能定理
\begin{tikzpicture}[scale=0.9,>=Latex]
  \draw[thick] (-1,0) -- (7,0);
  % 起点
  \draw[thick] (0,0) rectangle (1.5,1.2);
  \node at (0.75,0.6) {$m$};
  \draw[->,thick,blue] (0.20,1.45) -- (1.30,1.45) node[midway,above] {$v_1$};
  % 终点
  \draw[thick] (5,0) rectangle (6.5,1.2);
  \node at (5.75,0.6) {$m$};
  \draw[->,thick,blue] (5.20,1.45) -- (6.30,1.45) node[midway,above] {$v_2$};
  % 恒定合力画在物块上，避免误看成连接两个状态的连线。
  \draw[->,ultra thick,red] (0.75,0.6) -- (2.25,0.6)
    node[midway,above] {$F_{\text{合}}$};
  \draw[->,thick,gray] (2.55,0.6) -- (4.45,0.6)
    node[midway,above,font=\small] {运动过程};
  % 位移
  \draw[<->] (0,-0.6) -- (6.5,-0.6) node[midway,below] {$x$};
  % 动能定理
  \node at (3.25,2.20) {$F_{\text{合}} x = \frac12 m v_2^2 - \frac12 m v_1^2$};
\end{tikzpicture}
